\section{What Approval Does and Does Not Oblige}\label{sec:obliges}

\subsection{The \term{fides humana} boundary}

The doctrinal architecture is fixed and simple. Public Revelation is
complete in Christ: ``In Christ, God has said everything, that is, he
has revealed himself completely,'' as the 2000 commentary puts it,
``and therefore Revelation came to an end with the fulfilment of the
mystery of Christ.'' Recognized private revelations do not belong to
the deposit of faith; in the Catechism's words, ``it is not their
role to complete Christ's definitive Revelation, but to help live
more fully by it in a certain period of history'' (CCC 67). The assent owed to Revelation is divine and Catholic
faith; the assent invited by an approved apparition is, in
Lambertini's rule as retrieved by the 2000 commentary, ``an assent of
human faith in keeping with the requirements of prudence''---and an
assent of Catholic faith to it ``is not even possible.'' Approval
means the message contains nothing against faith or morals, may be
published, and may prudently be believed; ``it is a help which is
offered, but which one is not obliged to use.''\endnote{CDF,
\work{Message of Fatima}, theological commentary, whose text supplies
both quotations (the completeness formulation, and CCC 67 as quoted
there); Lambertini and Dhanis as in \S\ref{sec:interp2000}. The 2024 norms restate
the same boundary for the current discipline: even a \term{Nihil
obstat} authorizes only prudent adherence and never certifies
supernatural origin (nn.~17, 23; art.~22 \S2).}

Concretely, for F\'atima: a Catholic is free---reasonably,
prudently, without disloyalty---to believe that Our Lady appeared at
the Cova da Iria, and free not to. What no Catholic is free to do is
invert the order of assent: to treat the message as a second deposit
of faith, an oracle that settles doctrine, binds consciences,
predicts obligatory futures, directs political allegiance, or
substitutes for Scripture, liturgy, and the sacraments. The Church's
whole handling of the case---the bishop's eight years of reserve, the
restrained 1930 formula, forty-three years of unpublished custody of
the third part, the 2000 commentary's insistence that prophecy
``show[s] the right path to take for the future'' rather than
satisfying curiosity---is one long enforcement of that order.

The 2024 norms put the same boundary in current procedural language,
and their formulation is worth having exactly, because it is the rule
now in force: when the Dicastery grants a \term{Nihil obstat},
``such phenomena do not become objects of faith, which means the
faithful are not obliged to give an assent of faith to them.
Rather, as in the case of charisms recognized by the Church, they are
`ways to deepen one's knowledge of Christ and to give oneself more
generously to him, while rooting oneself more and more deeply in
communion with the entire Christian people'.'' The norms' own
presentation names F\'atima as the standing illustration of what
approval means.\endnote{DDF, \work{Norms for Proceeding in the
Discernment of Alleged Supernatural Phenomena} (17 May 2024), n.~12,
quoted exactly (the inner quotation is the norms' own, at their note
15); presentation, on the 2000 F\'atima explanation. The 1930 act is
a legacy judgment and is not relabeled under this taxonomy
(\S\ref{sec:standing}).}

\subsection{Against fatalism}

Everything in the F\'atima corpus that most attracts a certain kind
of reader---the punishments, the annihilated nations, the martyred
good---is precisely what the Church's own commentary reads as
conditional. The commentary is unusually blunt about it, and its
sentences are the sharpest instrument the tradition has against
apocalyptic determinism.

``The future is not in fact unchangeably set, and the image which the
children saw is in no way a film preview of a future in which nothing
can be changed.'' ``The purpose of the vision is not to show a film
of an irrevocably fixed future. Its meaning is exactly the opposite:
it is meant to mobilize the forces of change in the right
direction.'' And with a named application: ``Therefore we must
totally discount fatalistic explanations of the `secret', such as,
for example, the claim that the would-be assassin of 13 May 1981 was
merely an instrument of the divine plan guided by Providence and
could not therefore have acted freely, or other similar ideas in
circulation.''\endnote{CDF, \work{Message of Fatima}, theological
commentary, ``An attempt to interpret the `secret' of Fatima,''
quoted exactly.}

This is not a devotional softening of a hard text. It is the
classical doctrine of conditional prophecy, which Aquinas states as a
division of prophecy itself: where the future is known ``as in its
cause,'' the resulting prophecy of denunciation ``is not always
fulfilled,'' since it announces a causal tendency that may be
``hindered by some other occurrence supervening''
(\S\ref{sec:mode}). The second part of the F\'atima secret is
conditional on its face---``\emph{if} people do not cease offending
God''; ``\emph{if} my requests are heeded''---and Sister Lucia read
it that way herself in 1982: ``it is people themselves who are
preparing their own punishment\ldots\ God warns us and calls us to
the right path, while respecting the freedom he has given us; hence
people are responsible.''

Three practical consequences follow, and they exclude most of what
circulates under F\'atima's name.

\begin{itemize}
\item \textbf{No date-setting.} A conditional announcement cannot
yield a calendar. Any construction that assigns a year to the
``period of peace,'' the chastisement, or the triumph has converted a
warning into a schedule and has thereby contradicted the only
authoritative interpretation the text has received.
\item \textbf{No decoding.} The rule that ``not every element of the
vision has to have a specific historical sense'' forbids treating the
city, the soldiers, the arrows, or the ``etc.'' of the Fourth Memoir
as a cipher for named persons or events (\S\ref{sec:mode},
\S\ref{sec:withheld}).
\item \textbf{No weaponizing.} A private revelation whose criterion
of truth is its orientation to Christ cannot be turned into an
instrument for impugning the legitimacy of popes, councils, or
episcopal acts. The commentary's own test---a message that ``becomes
independent of him or even presents itself as another and better plan
of salvation, more important than the Gospel\ldots\ certainly does
not come from the Holy Spirit''---applies to interpretations as much
as to messages.
\end{itemize}

A current safeguard of the same kind, from the same dicastery,
belongs here: the doctrinal note \work{Mater Populi fidelis} of
4 November 2025 opens by observing that certain Marian titles ``are
frequently used in connection with'' alleged supernatural phenomena
and ``are not always employed precisely,'' with consequences that
``can often lead to a mistaken understanding of Mary's role.''
Nothing in the F\'atima texts printed in \S\ref{sec:text} attributes
to Mary an independent power; the promise ``my Immaculate Heart will
triumph'' is read by the 2000 commentary through her \term{fiat},
which brought the Saviour into the world, and is not to be read as a
rival economy of salvation.\endnote{DDF, \work{Mater Populi fidelis}
(4 November 2025), n.~2, English text at
\url{https://www.vatican.va/roman_curia/congregations/cfaith/documents/rc_ddf_doc_20251104_mater-populi-fidelis_en.html}
(read 2026-07-25), quoted exactly; the note's own footnote at that
place identifies the titles concerned. Cited here for its bearing on
apparition-linked Marian language generally; it does not name
F\'atima. CDF, \work{Message of Fatima}, theological commentary, for
the reading of the triumph through Mary's \term{fiat}.}

\subsection{What the positive judgments leave unjudged}

Each act reaches its own object, and the gaps are part of the record.
The 1930 declaration does not reach the angel apparitions of 1916,
the Pontevedra and Tuy experiences, or any text of the memoirs and
the secret---these stand on the Holy See's later practical reception,
which published, interpreted, and built devotion upon them without
ever issuing a formal supernaturality judgment on them. The
canonizations certify the sanctity of Francisco and Jacinta, not the
accuracy of every remembered dialogue. The consecrations are acts of
papal piety defined by their texts. The 2000 publication closes the
textual question of the secret and gives the Church's interpretation;
its own author calls the interpretation ``an attempt.'' And no act,
diocesan or Roman, has ever judged the mechanism of the solar
phenomenon, guaranteed the verbal inerrancy of the children's
reported sentences, or authenticated the ``devotional accretions,
promises, and private interpretations'' that circulate under
F\'atima's name. Where a claim exceeds its act, the excess is the
claimant's.\endnote{For the reported-message discipline observed
throughout this study: a sentence is the Lady's only in the sense
that a named witness reported it at a stated date; the phrase
``reported message'' is the study's default register, and the
strata (1917 records; 1924 sworn interrogation; 1935--1941 memoirs;
1944 manuscript) are identified at each use.}

\subsection{The message's evangelical center}

Read under those controls, what the approved F\'atima proposes is
almost severe in its simplicity: conversion; penance; the daily
Rosary (asked, in the 1917 records, at every apparition); reparation,
above all Eucharistic, for sin; intercession for sinners---``many
souls go to hell, because there are none to sacrifice themselves and
to pray for them,'' in the memoirs' account of the Valinhos words;
devotion to the Immaculate Heart as schooling in Mary's
\term{fiat}; and peace in history as a fruit hanging on human
freedom before God. Everything in it leads back into the Gospel it
cannot supplement---which is, by the Church's stated criterion, the
very mark of its credibility.\endnote{Center as stated by the 2000
commentary (``the exhortation to prayer as the path of `salvation
for souls' and, likewise, the summons to penance and conversion'';
Lucia to Ratzinger: the purpose of the apparitions was growth in
faith, hope, and love, ``everything else was intended to lead to
this''); Rosary requests in the 1917 records
(\work{Sele\c{c}\~ao}, Docs.~1, 3, 14); sacrifice-for-sinners
wording from the memoir stratum as printed by the Sanctuary's
narrative page.}
